\section*{第 \ref{ch:g12:human-evolution} 章——人类演化一瞥}

\begin{solution}{exo:g12:human-evolution:1}
黑猩猩（与倭黑猩猩一起）；两条谱系
在约六到七百万年前分开。
\end{solution}

\begin{solution}{exo:g12:human-evolution:2}
颅骨下方的脊柱开口；双弯曲的脊柱；
短而呈碗状的骨盆；股骨向内倾斜至膝部；
足呈拱形，大脚趾与其他脚趾排成一线。
\end{solution}

\begin{solution}{exo:g12:human-evolution:3}
约 \qty{420}{cm^3}；约 \qty{850}{cm^3}；约 \qty{1350}{cm^3}。
\end{solution}

\begin{solution}{exo:g12:human-evolution:4}
由于谱系中有多个\omterm{def:g10:biodiversity-scales:species}{物种}同时存在，其中大多数
最终灭绝，而非演化为下一个\omterm{def:g10:biodiversity-scales:species}{物种}；
\omterm{def:g10:biodiversity-scales:species}{物种}是分支，而非阶梯。
\end{solution}

\begin{solution}{exo:g12:human-evolution:5}
约 300\,000 年前，发生于非洲。\omterm{def:g10:biodiversity-scales:biodiversity}{遗传多样性}在非洲最高，
并随距非洲距离增加而降低：所有现存人类
均源自这一非洲\omterm{def:g12:selection-drift-speciation:population}{种群}，
经由一系列相继出现的小型奠基群体扩散而来。
\end{solution}

\begin{solution}{exo:g12:human-evolution:6}
\omterm{def:g12:human-evolution:lineage}{人类谱系}中两条祖先\omterm{def:g11:cell-cycle-mitosis:chromosome}{染色体}融合为一条，因此
24 对变为 23 对。融合后的\omterm{def:g11:cell-cycle-mitosis:chromosome}{染色体}中部保留了两条原始末端的残迹，
以及第二个沉默的\omterm{def:g11:cell-cycle-mitosis:chromosome}{着丝粒}——这是融合的记录；
其\omterm{def:g10:universal-dna:gene}{基因}内容与两条黑猩猩\omterm{def:g11:cell-cycle-mitosis:chromosome}{染色体}逐一对应。
源自一套共同的\omterm{def:g11:cell-cycle-mitosis:chromosome}{染色体}，
仅经一次修改，正是亲缘关系所预测的。
\end{solution}

\begin{solution}{exo:g12:human-evolution:7}
直立行走，约在二百万年前出现：足迹和骨盆
距今 3.6 至 3.2 百万年，属于脑容量为 \qty{400}{cm^3} 的行走者；
脑的增长始于约 2.5 百万年前。
\end{solution}

\begin{solution}{exo:g12:human-evolution:8}
二百万年前：\emph{A. africanus}（勉强）、\emph{Paranthropus}、
\emph{Homo habilis} 和早期的 \emph{H. erectus}。约 100\,000 年前：
\emph{H. erectus}（最后的\omterm{def:g12:selection-drift-speciation:population}{种群}）、\emph{H. heidelbergensis}（最后）、
尼安德特人和 \emph{H. sapiens}。
\end{solution}

\begin{solution}{exo:g12:human-evolution:9}
走出非洲的现代人与尼安德特人之间的杂交，
发生于约 50\,000 年前的西亚和欧洲。留在撒哈拉以南非洲的\omterm{def:g12:selection-drift-speciation:population}{种群}
从未与尼安德特人相遇，因此
不携带任何尼安德特人的\omterm{def:g10:universal-dna:gene}{基因}。
\end{solution}

\begin{solution}{exo:g12:human-evolution:10}
直立行走（枕骨大孔靠前）、脑小、颌骨粗大，约 2.8 百万年前：
一种南方古猿，被置于谱系基部附近的一条侧支上——
是 \emph{Homo} 祖先的近亲，
而非已被证实的祖先。
\end{solution}

\begin{solution}{exo:g12:human-evolution:11}
大猿类构成一个\omterm{prop:g12:phylogenetic-trees:rules}{演化支}，人类是其中的一个分支；
由于\omterm{prop:g12:phylogenetic-trees:rules}{演化支}包含其节点的全部后代，人类也是猿，
正如人类同样是灵长类和哺乳动物。“源自猿类”的说法错误地把
猿类放在一个节点上；现存的猿类是互为姐妹的末端类群。
\end{solution}

\begin{solution}{exo:g12:human-evolution:12}
假说 1：竞争——现代人数量更多或组织更好，
夺取了资源；证据将是两者重叠区域内尼安德特人遗址的减少。假说 2：吸收
——尼安德特人小\omterm{def:g12:selection-drift-speciation:population}{种群}部分通过杂交被吸收，
部分则因\omterm{prop:g12:selection-drift-speciation:drift}{遗传漂变}和偶然性，在气候变化中灭绝；
证据将是人类体内的尼安德特人 \omterm{def:g10:universal-dna:information}{DNA}，以及接触前
小而衰退的尼安德特人\omterm{def:g12:selection-drift-speciation:population}{种群}的迹象。
两者可能都成立。
\end{solution}

\begin{solution}{exo:g12:human-evolution:13}
扩散的每一步都由一个小群体完成，该群体携带其所离开\omterm{def:g12:selection-drift-speciation:population}{种群}的
部分\omterm{def:g10:universal-dna:gene}{等位基因}样本；每一个新区域都由前一区域奠基而成。
每次奠基多样性都会下降，因此最远的\omterm{def:g12:selection-drift-speciation:population}{种群}多样性最低。
扩散是通过相继发生的小规模迁移完成的，
而非大规模的整体迁移。
\end{solution}

\begin{solution}{exo:g12:human-evolution:14}
畜牧业使乳汁成为成年人的食物：携带乳糖酶持续表达的\omterm{def:g10:universal-dna:gene}{等位基因}者
得以更好地进食，并留下更多后代。迁移到日照微弱的纬度，
使维生素 D 变得稀缺：肤色较浅、能合成更多维生素 D 者
更为有利。农业造成人口密集并形成积水，疟疾由此传播：
能抵御疟疾的\omterm{def:g10:universal-dna:gene}{等位基因}随之上升。在每种情况中，
生活方式的改变都改变了哪些\omterm{def:g10:universal-dna:gene}{等位基因}得以传递。
\end{solution}

\begin{solution}{exo:g12:human-evolution:15}
选择是指携带不同\omterm{def:g10:universal-dna:gene}{等位基因}者之间繁殖上的任何差异；
医学改变的是哪些\omterm{def:g10:universal-dna:gene}{等位基因}重要，而非是否有\omterm{def:g10:universal-dna:gene}{等位基因}重要——对感染的抗性、
生育力、生殖年龄仍有差异，且可遗传。这三种性状表明，
\omterm{def:g12:selection-drift-speciation:selection}{自然选择}在过去一万年内持续作用，
由文化本身所驱动。与此同时，漂变在每个\omterm{def:g12:selection-drift-speciation:population}{种群}中
始终存在，无论如何。人类演化已改变方向，
而非停止。
\end{solution}

\begin{solution}{pb:g12:human-evolution:1}
\textbf{1.} A：$21.5/0.15 \approx \qty{1.43}{m}$；B：$18.5/0.15
\approx \qty{1.23}{m}$。

\textbf{2.} $0.8 \times 1.43 = \qty{1.14}{m}$，$0.8 \times 1.23 =
\qty{0.98}{m}$：两者步幅均短于慢走，说明他们当时
是从容行走。

\textbf{3.} 足跟深凹：足跟先着地，与人类行走时相同；
足弓：足部起杠杆和弹簧的作用；大脚趾与其他趾平行：
足部用于蹬地，而非抓握树枝。

\textbf{4.} 足跟先着地、步幅较大的步态，足部类人，
持续行进 \qty{27}{m}：习惯性双足行走。仅凭足迹无法得知：
骨盆、股骨角度、脊柱开口及脊柱的弯曲。

\textbf{5.} 直立行走先于大脑的扩大
出现，且早了一百多万年。

\textbf{6.} 1：南方古猿；2：\emph{Homo habilis}；3：\emph{Homo
erectus}；4：尼安德特人；5：\emph{Homo sapiens}。

\textbf{7.} 由 3.0 百万年前的 450 增至 1.9 百万年前的 650，1.0 百万年前为 950，随后
为 1450 和 1350：增长最快的是 1.9 至 1.0 百万年前之间
（每百万年 \qty{300}{cm^3}），其次是 1.0 至 0.06 百万年前之间。

\textbf{8.} 不是：尼安德特人的脑与之一样大，甚至更大。现代人的鉴别特征是
额头高、脑颅下方的面部较小、
有下巴以及圆形的颅骨。

\textbf{9.} 演化丛中并存的两个类群，解剖结构不同，
却能相互杂交：现今生活的人群中，
非洲以外人群的尼安德特人 \omterm{def:g10:universal-dna:information}{DNA} 即是证据。

\textbf{10.} 化石只是其自身侧枝上的一个末端；其性状组合
表明它靠近后代\omterm{def:g10:biodiversity-scales:species}{物种}所由分出的节点，
但没有任何性状能证明它就是产生这些后代的那个\omterm{def:g12:selection-drift-speciation:population}{种群}。
它是祖先的近亲，
这正是化石所能证明的全部。

\textbf{11.} 构成她\omterm{def:g10:universal-dna:gene}{基因}组 1.6\% 的片段，
与尼安德特人的序列比与任何现代人的序列更为接近：
它们遗传自与现代人杂交的尼安德特人祖先。

\textbf{12.} 图表显示，欧洲人群中尼安德特人成分约为 1.8\%，
美拉尼西亚的亚洲人群的 \omterm{def:g10:universal-dna:information}{DNA} 中约为 3.5\%：她的欧洲血统带来
尼安德特人的份额，她的美拉尼西亚血统则带来另一部分
（因其混合血统而被稀释）。

\textbf{13.} $50\,000/25 = 2000$ 代。该份额未被稀释，
因为整个\omterm{def:g12:selection-drift-speciation:population}{种群}都携带它：与其他携带者混合，
平均值保持不变；只有携带者与非携带者混合时才会被稀释，
而在非洲以外地区并不存在非携带者。

\textbf{14.} 这说明两者之间的界限并未封闭：分离约 500\,000
年后，两个\omterm{def:g12:selection-drift-speciation:population}{种群}仍能产生可育后代，
后代——正在形成中的\omterm{def:g10:biodiversity-scales:species}{物种}，与其说是被生物学所隔开，
不如说更多是被地理所隔开。

\textbf{15.} \omterm{def:g10:cells-common-unit:organelle}{线粒体}只经母亲遗传：现今人类中
未保留尼安德特人的母系谱系，无论是因为尼安德特人母亲的后代
已经灭绝，还是因为杂交大多只朝一个方向进行。
核\omterm{def:g10:universal-dna:gene}{基因}中的痕迹得以保留，而该母系谱系则未能保留。

\textbf{16.} $1.2/0.27 \approx 4.4$ 百万年——偏低的估计，
因为分子钟是依据更远的节点校准的。

\textbf{17.} 化石比分子钟的估计更古老；真实的
节点至少与该谱系最古老的化石一样古老，即 6 至 7 百万年，
且分子钟在这一较近时段的速率必定略慢于
0.27\%。

\textbf{18.} \emph{H. erectus}，以及时间线上最后的
\emph{Paranthropus}（以及自约 0.7 百万年前起的 \emph{H. heidelbergensis}）；
在这些类群中，只有经由 \emph{H. heidelbergensis}
通向 \emph{H. sapiens} 的谱系有现存后代，而 \emph{H. erectus}
仅通过汇入它的\omterm{def:g12:selection-drift-speciation:population}{种群}延续下来——这只是若干支系中的一支。

\textbf{19.} 在这段时间里，骨骼和大脑变化甚微，
而工具、艺术、语言、农业和城市则改变了人类的生活方式；
这些变化通过学习传递，在一代之内即可积累，
并且如今正在塑造\omterm{def:g10:biodiversity-scales:species}{物种}演化所处的
环境。

\textbf{20.} 沃克 A 身高约 \qty{1.4}{m}，行走缓慢，脚与
类似人类的足部；足部先出现，大脑晚了一百万年，
下巴最后出现；这位学生携带着两个已灭绝\omterm{def:g12:selection-drift-speciation:population}{种群}的\omterm{def:g10:universal-dna:gene}{基因}。
\end{solution}
